% !TEX spellcheck = en_US
% !TEX spellcheck = LaTeX
\documentclass[letterpaper,10pt,english]{article}
\usepackage{%
amsfonts,%
amsmath,%
amssymb,%
amsthm,%
babel,%
bbm,%
%biblatex,%
caption,%
centernot,%
color,%
enumerate,%
%enumitem,%
epsfig,%
epstopdf,%
etex,%
fancybox,%
framed,%
fullpage,%
%geometry,%
graphicx,%
hyperref,%
latexsym,%
mathptmx,%
mathtools,%
multicol,%
paralist,%
pgf,%
pgfplots,%
pgfplotstable,%
pgfpages,%
proof,%
psfrag,%
%subfigure,%
tcolorbox,%
tfrupee,%
tikz,%
times,%
%ulem,%
url,%
xcolor,%
mathpazo,%
}

\definecolor{shadecolor}{gray}{.95}%{rgb}{1,0,0}
\usepackage[margin=1in,top=0.75in]{geometry}
\usepackage[mathscr]{eucal}
\usepgflibrary{shapes}
\usepgfplotslibrary{fillbetween}
\usetikzlibrary{%
arrows,%
backgrounds,%
chains,%
decorations.pathmorphing,% /pgf/decoration/random steps | erste Graphik
decorations.text,% 
matrix,%
positioning,% wg. " of "
fit,%
patterns,%
petri,%
plotmarks,%
scopes,%
shadows,%
shapes.misc,% wg. rounded rectangle
shapes.arrows,%
shapes.callouts,%
shapes%
}

%\pgfplotsset{compat=newest} %<------ Here
\pgfplotsset{compat=1.11} %<------ Or use this one

\theoremstyle{plain}
\newtheorem{thm}{Theorem}[section]
\newtheorem{lem}[thm]{Lemma}
\newtheorem{prop}[thm]{Proposition}
\newtheorem{cor}[thm]{Corollary}
\newtheorem{clm}[thm]{Claim}

\theoremstyle{definition}
\newtheorem{defn}[thm]{Definition}
\newtheorem{conj}[thm]{Conjecture}
\newtheorem{exmp}[thm]{Example}
\newtheorem{axiom}[thm]{Axiom}
\newtheorem{assum}[thm]{Assumption}
\newtheorem{exerc}[thm]{Exercise}
\newtheorem*{defin}{Definition}

\theoremstyle{remark}
\newtheorem{rem}{Remark}
\newtheorem{note}{Note}

\newcommand{\iid}{\emph{i.i.d.}\ }
\newcommand{\Cov}{\operatorname{Cov}}
%\newcommand{\det}{\operatorname{det}}
%\newcommand{\Real}{\mathbb{R}}
\newcommand{\tr}{\operatorname{tr}}
\newcommand{\Var}{\operatorname{Var}}
\newcommand{\cov}{\operatorname{cov}}

%\renewcommand{\proof}[1]{\begin{proof}#1\end{proof}}
\newcommand{\EQ}[1]{\begin{equation*}#1\end{equation*}}
\newcommand{\EQN}[1]{\begin{equation}#1\end{equation}}
\newcommand{\eq}[1]{\begin{align*}#1\end{align*}}
\newcommand{\meq}[2]{\begin{xalignat*}{#1}#2\end{xalignat*}}
\newcommand{\norm}[1]{\left\lVert#1\right\rVert}
\newcommand{\inner}[1]{\left\langle #1\right\rangle}
\newcommand{\abs}[1]{\left\lvert#1\right\rvert}
\newcommand{\expect}[1]{\mathbb{E}\left[{#1}\right]}
\newcommand{\prob}[1]{\mathbb{P}\left[{#1}\right]}
\newcommand{\given}{\; \big\vert \;} 
\newcommand{\set}[1]{\left\{#1\right\}} 
\newcommand{\indicator}[1]{1_{\set{#1}}} 
\newcommand{\Ind}[1]{\mathbbm{1}_{#1}} 
\newcommand{\SetIn}[1]{\mathbbm{1}_{\set{#1}}} 
\newcommand{\red}[1]{\textcolor{red}{#1}} 
\newcommand{\floor}[1]{\left\lfloor#1\right\rfloor}
\newcommand{\ceil}[1]{\left\lceil#1\right\rceil}


\newcommand{\C}{\mathbb{C}}
\newcommand{\D}{\mathbb{D}}
\newcommand{\E}{\mathbb{E}}
\newcommand{\F}{\mathbb{F}}
\newcommand{\N}{\mathbb{N}}
\renewcommand{\P}{\mathbb{P}}
\newcommand{\Q}{\mathbb{Q}}
\newcommand{\R}{\mathbb{R}}
\newcommand{\Z}{\mathbb{Z}}

\newcommand{\bA}{\mathbf{A}}
\newcommand{\bI}{\mathbf{I}}
\newcommand{\bU}{\mathbf{1}}
\newcommand{\bx}{\mathbf{x}}
\newcommand{\bX}{\mathbf{X}}
\newcommand{\bZ}{\mathbf{Z}}


\newcommand{\cA}{\mathcal{A}}
\newcommand{\cB}{\mathcal{B}}
\newcommand{\cC}{\mathcal{C}}
\newcommand{\cF}{\mathcal{F}}
\newcommand{\cG}{\mathcal{G}}
\newcommand{\cM}{\mathcal{M}}
\newcommand{\cN}{\mathcal{N}}
\newcommand{\cP}{\mathcal{P}}
\newcommand{\cT}{\mathcal{T}}
\newcommand{\cX}{\mathcal{X}}
\newcommand{\cY}{\mathcal{Y}}

\newcommand{\sA}{\mathscr{A}}
\newcommand{\sB}{\mathscr{B}}
\newcommand{\sC}{\mathscr{C}}
\newcommand{\sE}{\mathscr{E}}
\newcommand{\sF}{\mathscr{F}}
\newcommand{\sG}{\mathscr{G}}
\newcommand{\sH}{\mathscr{H}}
\newcommand{\sL}{\mathscr{L}}
\newcommand{\sS}{\mathscr{S}}
\newcommand{\sT}{\mathscr{T}}
\newcommand{\sX}{\mathscr{X}}
\newcommand{\sY}{\mathscr{Y}}
\newcommand{\sZ}{\mathscr{Z}}

\newcommand{\tS}{\tilde{S}}
% Debug
\newcommand{\todo}[1]{\begin{color}{blue}{{\bf~[TODO:~#1]}}\end{color}}

% a few handy macros

\renewcommand{\le}{\leqslant}
\renewcommand{\ge}{\geqslant}
\newcommand\matlab{{\sc matlab}}
\newcommand{\goto}{\rightarrow}
\newcommand{\bigo}{{\mathcal O}}
%\newcommand{\half}{\frac{1}{2}}
%\newcommand\implies{\quad\Longrightarrow\quad}
\newcommand\reals{{{\rm l} \kern -.15em {\rm R} }}
\newcommand\complex{{\raisebox{.043ex}{\rule{0.07em}{1.56ex}} \hskip -.35em {\rm C}}}


% macros for matrices/vectors:

% matrix environment for vectors or matrices where elements are centered
\newenvironment{mat}{\left[\begin{array}{ccccccccccccccc}}{\end{array}\right]}
\newcommand\bcm{\begin{mat}}
\newcommand\ecm{\end{mat}}

% matrix environment for vectors or matrices where elements are right justifvied
\newenvironment{rmat}{\left[\begin{array}{rrrrrrrrrrrrr}}{\end{array}\right]}
\newcommand\brm{\begin{rmat}}
\newcommand\erm{\end{rmat}}

% for left brace and a set of choices
%\newenvironment{choices}{\left\{ \begin{array}{ll}}{\end{array}\right.}
\newcommand\when{&\text{if~}}
\newcommand\otherwise{&\text{otherwise}}
% sample usage:
%\delta_{ij} = \begin{choices} 1 \when i=j, \\ 0 \otherwise \end{choices}


% for labeling and referencing equations:
\newcommand{\eql}{\begin{equation}\label}
\newcommand{\eqn}[1]{(\ref{#1})}
% can then do
%\eql{eqnlabel}
%...
%\end{equation}
% and refer to it as equation \eqn{eqnlabel}.
% some useful macros for finite difference methods:
\newcommand\unp{U^{n+1}}
\newcommand\unm{U^{n-1}}
% for chemical reactions:
\newcommand{\react}[1]{\stackrel{K_{#1}}{\rightarrow}}
\newcommand{\reactb}[2]{\stackrel{K_{#1}}{~\stackrel{\rightleftharpoons}
 {\scriptstyle K_{#2}}}~}
\makeatletter
\def\th@plain{%
\thm@notefont{}% same as heading font
\itshape % body font
}
\def\th@definition{%
\thm@notefont{}% same as heading font
\normalfont % body font
}
\makeatother
\date{}


\title{Lecture-12: Conditional Expectation}
\author{}

\begin{document}
\maketitle

\section{Conditional expectation given a non trivial event}
Consider a probability space $(\Omega, \sF, P)$ and an event $B \in \sF$ such that $P(B) > 0$. 
Then, the conditional probability of any event $A \in \sF$ given an event $B$ is defined as
\EQ{
P(A \given B) = \frac{P(A\cap B)}{P(B)}.
}
Consider a random variable $X: \Omega\to\R$ defined on a probability space $(\Omega, \sF, P)$, 
with distribution function $F_X: \R \to [0,1]$, 
and a non trivial event $B \in \sF$ such that $P(B) > 0$. 
\begin{defn}%[Conditional distribution given a non trivial event] 
%Consider a random variable $X: \Omega\to\R$ defined on a probability space $(\Omega, \sF, P)$, 
%and a non trivial event $B \in \sF$ such that $P(B) > 0$, then 
The \textbf{conditional distribution of $X$ given event $B$} is denoted by $F_{X| B}: \R \to [0,1]$ 
and $F_{X| B}(x)$ is defined as the probability of event $A_X(x) \triangleq X^{-1}(-\infty, x]$ conditioned on event $B$ for all $x\in\R$. 
That is, 
\EQ{
F_{X| B}(x) \triangleq P(A_X(x)| B) = \frac{P(A_X(x)\cap B)}{P(B)}\text{ for all }x \in \R. 
}
\end{defn}
\begin{rem} 
%For a random variable $X:\Omega \to \R$ and a non trivial event $B \in \sF$ defined on a probability space $(\Omega,\sF, P)$, 
The conditional distribution $F_{X|B}$ is a distribution function. 
This follows from the fact that 
%\begin{compactenum}[(i)]
%\item 
(i) $F_{X|B} \ge 0$,
%\item 
(ii) $F_{X|B}$ is right continuous,  
%\item 
(iii) $\lim_{x\downarrow -\infty}F_{X|B}(x) = 0$ and $\lim_{x\uparrow \infty}F_{X|B}(x) = 1$.
%\end{compactenum}
\end{rem}
\begin{rem}
For a discrete random variable $X:\Omega\to\sX$, the conditional probability mass function of $X$ given a non trivial event $B$ is given by $P_{X|B}(x) = \frac{P(X^{-1}\set{x}\cap B)}{P(B)}$ for all $x \in \sX$. 
\end{rem}
\begin{rem}
For a continuous random variable $X:\Omega\to\R$, the conditional density of $X$ given a non trivial event $B$ is given by $f_{X|B}(x) = \frac{dF_{X|B}(x)}{dx}$ for all $x \in \R$. 
\end{rem}
\begin{shaded}
\begin{exmp}[Conditional distribution] 
\label{exmp:FairDieConditionHalf}
Consider the probability space $(\Omega, \sF, P)$ corresponding to a random experiment where a fair die is rolled once. 
For this case, the outcome space $\Omega = [6]$, the event space $\sF = \cP([6])$, and the probability measure $P(\omega) = \frac{1}{6}$ for all $\omega\in \Omega$. 

We define a random variable $X: \Omega\to\R$ such that $X(\omega) = \omega$ for all $\omega \in \Omega$, 
and an event $B \triangleq \set{\omega \in \Omega: X(\omega) \le 3} = [3] \in \sF$. 
We note that $P(B) = 0.5$ and the conditional %distribution 
PMF of $X$ given $B$ is 
\EQ{
%F_{X|B}(x) = \Ind{[3, \infty)}(x) + \frac 2 3\Ind{[2,3)}(x) + \frac 1 3\Ind{[1, 2)}(x).
P_{X|B}(x) =  \frac 1 3\SetIn{x=3} + \frac 1 3\SetIn{x=2} + \frac 1 3\SetIn{x=1}.
}
\end{exmp}
\end{shaded}

\begin{defn}
%Consider a random variable $X: \Omega\to\R$ defined on a probability space $(\Omega, \sF, P)$, 
%and a non trivial event $B \in \sF$ such that $P(B) > 0$, then 
The \textbf{conditional expectation of $X$ given event $B$} is given as 
$
\E[X|B] \triangleq \int_{x\in\R}xdF_{X|B}(x).
$
\end{defn}
\begin{rem}
For a discrete random variable $X:\Omega\to\sX$, the conditional expectation of $X$ given a non trivial event $B$ is given by $\E[X|B]= \sum_{x\in\sX}xP_{X|B}(x)$. 
\end{rem}
\begin{shaded}
\begin{exmp}[Conditional expectation]
For the random variable $X$ and event $B$ defined in Example~\ref{exmp:FairDieConditionHalf}, 
the conditional expectation $\E[X|B] = 2$. 
\end{exmp}
\end{shaded}
\begin{rem}
Consider two random variables $X,Y$ defined on this probability space, 
then for $y \in \R$ such that $F_Y(y) > 0$, we can define events $A_X(x) \triangleq X^{-1}(-\infty, x]$ and $A_Y(y) = Y^{-1}(-\infty, y]$, 
such that 
\EQ{
P\left(\set{X \le x}\given\set{Y \le y}\right) = \frac{F_{X,Y}(x,y)}{F_Y(y)}.
}
The key observation is that $\set{Y \le y}$ is a non-trivial event. 
How do we define conditional expectation based on events such as $\set{Y = y}$? 
When random variable $Y$ is continuous, this event has zero probability measure. 
\end{rem}

\section{Conditional expectation given an event space}
Consider random variables $X:\Omega\to\R$ and $Y:\Omega\to\R$ defined on the same probability space $(\Omega, \sF, P)$ 
such that $\E\abs{X} < \infty$, and a smaller event space $\sG \subset \sF$. 
For each non trivial event $G \in \sG$, we know how to define the conditional distribution $F_{X|G}$ and $\E[X|G]$. 
For any trivial event $N \in \sG$, these are undefined. 

%, %and we can define them as zero. 
\begin{defn}%[Conditional expectation given an event space]
%Consider a random variable $X:\Omega\to\R$ defined on a probability space $(\Omega, \sF, P)$ such that $\E\abs{X} < \infty$ and a smaller event space $\sG \subset \sF$. 
The \textbf{conditional expectation of random variable $X$ given event space $\sG$} is a random variable $\E[X\given \sG]: \Omega \to \R$ defined on the same probability space, 
such that 
\begin{compactenum}[(i)]
\item $Z \triangleq \E[X\given \sG]$ is $\sG$ measurable, %i.e. $Y^{-1}(-\infty, y] \in \sG$ for all $y \in \R$, 
\item for all $G\in \sG$, we have
$
%\int_AXdP = 
\E[X\Ind{G}] = \E[Z\Ind{G}], %=\int_AYdP,
$
\item $\E\abs{Z}<\infty$.
\end{compactenum}
\end{defn}
\begin{lem} 
%Consider a random variable $X:\Omega\to\R$ defined on a probability space $(\Omega,\sF, P)$ and a sub event space $\sG\subseteq \sF$. 
The conditional expectation of $X$ given $\sG$  is an a.s. unique random variable.
\end{lem}
\begin{proof}
Consider two random variables $Z_1 = \E[X|\sG]$ and $Z_2 = \E[X|\sG]$. 
Then from the definition, $Z_1, Z_2$ are $\sG$ measurable random variables, 
and $Z_1-Z_2$ is also $\sG$ measurable. 
Therefore, $G_n \triangleq \set{Z_1 - Z_2 > \frac{1}{n}} \in \sG$ and $\E[(Z_1-Z_2)\Ind{G_n}] = 0$ by definition. 
It follows from continuity of probability, that $P(\lim_n G_n) = 0$. 
Similarly, defining $F_n \triangleq \set{Z_2 - Z_1 > \frac 1 n}$, we can show that $P(\lim_nF_n) = 0$. 
%It follows that $P(\lim_n F_n^c\cap G_n^c) = 1$. 
\end{proof}
\begin{shaded}
\begin{exmp}[Conditional expectation as averaging]
Consider a random variable $X:\Omega\to\R$ defined on a probability space $(\Omega,\sF,P)$ with $\E\abs{X} < \infty$, 
and the coarsest event space $\sG = \set{\emptyset, \Omega} \subseteq \sF$ and finest event space $\sF$. 
We observe that $\E[X|\sG] = \E X$ a.s. uniquely, since 
%\begin{compactenum}[(i)]
%\item 
(i) $\E X$ is a constant and hence $\sG$ measurable, 
%\item 
(ii) $\E[\E X \Ind{\emptyset}] = \E[X \Ind{\emptyset}] = 0$ and $\E[\E X \Ind{\Omega}] = \E[X \Ind{\Omega}] = \E X$, and 
%\item 
(iii) $\E\abs{\E X} = \abs{\E X} \le \E \abs{X} < \infty$ from the Jensen's inequality. 
%\end{compactenum}

We also observe that $\E[X|\sF] = X$ a.s. uniquely, since 
%\begin{compactenum}[(i)]
%\item 
(i) $X$ is $\sF$ measurable random variable, 
%\item 
(ii) $\E[ X \Ind{G}] = \E[X \Ind{G}]$ for all events $G\in \sF$, and 
%\item 
(iii) $\E\abs{X} < \infty$. 
%\end{compactenum}
\end{exmp}
\end{shaded}
%\begin{rem}
%Expectation is averaging, and conditional expectation is averaging over event spaces. 
%%Consider a random variable $X:\Omega\to\R$ defined on a probability space $(\Omega,\sF,P)$ with $\E\abs{X} < \infty$.
%For the coarsest event space $\set{\emptyset, \Omega} \subseteq \sF$ and finest event space $\sF$, 
%we can observe that the coarsest averaging is $\E[X|\set{\emptyset, \Omega}] = \E X$ 
%and the finest averaging is $\E[X|\sF] = X$. 
%\end{rem}
\begin{lem} 
The mean of conditional expectation of random variable $X$ given event space $\sG$ is $\E X$. 
\end{lem}
\begin{proof} 
From the definition of event space $\Omega\in\sG$, 
and from the definition of conditional expectation, we get 
$
\E[\E[X\mid \sG]] 
= \E[\E[X\mid \sG]\Ind{\Omega}] 
= \E[X\Ind{\Omega}]
= \E X.  
$
\end{proof}
%\subsection{Conditional expectation conditioned on a random variable}
\begin{defn}%[Conditional expectation given a random variable]
%Consider two random variables $X, Y:\Omega\to\R$ defined on the same probability space $(\Omega, \sF, P)$ such that $\E\abs{X} < \infty$. 
The conditional expectation of $X$ given $Y$ is a random variable $\E[X\given Y] \triangleq \E[X\given\sigma(Y)]$ defined on the same probability space. 
\end{defn}
\begin{shaded}
\begin{exmp}[Conditioning on simple random variables] 
For a simple random variable $Y: \Omega \to \sY \subseteq \R$ defined on the probability space $(\Omega, \sF, P)$, 
we define fundamental events $E_y \triangleq Y^{-1}\set{y} \in \sF$ for all $y \in \sY$. 
Then the sequence of events $E \triangleq (E_y \in \sF: y \in \sY)$ partitions the sample space, 
and we can write the event space generated by random variable $Y$ as $\sigma(Y) = (\cup_{y\in I}E_y: I \subseteq \sY)$. 

For a random variable $X: \Omega \to \R$ defined on the same probability space, 
the random variable $Z \triangleq \E[X\given Y]$ is $\sigma(Y)$ measurable. 
Therefore, $\E[X| Y] = \sum_{y\in\sY}\alpha_y\Ind{E_y}$ for some $\alpha\in\R^\sY$. 
We verify that $Z:\Omega\to\R$ is a $\sigma(Y)$ measurable random variable, 
since $\sigma(Z) \subseteq \sigma(Y)$. 
We can also check that $\E\abs{Z} < \infty$. 
Further, we have $\E[Z\Ind{E_y}] = \E[X\Ind{E_y}]$ for any $y \in \sY$, 
which implies that $\alpha_y = \frac{\E[X\Ind{E_y}]}{P_Y(y)}$ for any $y \in \sY$.
Notice that 
\EQ{
\E[X\mid E_y] 
= \int_{x\in\R}x dF_{X\mid E_y}(x) 
= \frac{1}{P_Y(y)}\int_{x\in\R}x d P(A_X(x)\cap E_y)
= \frac{1}{P_Y(y)}\E[X\Ind{E_y}]
= \alpha_y.
}
\end{exmp}
\end{shaded}
\begin{rem} 
There are three main takeaways from this definition. 
For a random variable $Y$, 
the event space generated by $Y$ is $\sigma(Y)$.  
\begin{enumerate}
\item 
The conditional expectation $\E[X|Y] = \E[X|\sigma(Y)]$ and is $\sigma(Y)$ measurable. 
That is, $\E[X|Y]$ is a Borel measurable function of $Y$. 
In particular when $Y$ is discrete, this implies that $\E[X|Y]$ is a simple random variable that takes value $\E[X|E_y]$ when $\omega \in E_y$, and the probability of this event is $P_Y(y)$. 
When $Y$ is continuous, $\E[X|Y]$ is a continuous random variable with density $f_Y$. 

\item Expectation is averaging. 
Conditional expectation is averaging over event spaces. 
We can observe that the coarsest averaging is $\E[X|\set{\emptyset, \Omega}] = \E X$ 
and the finest averaging is $\E[X|\sigma(X)] = X$. 
Further, $\E[X|\sigma(Y)]$ is averaging of $X$ over events generated by $Y$.   
If we take any event $A \in \sigma(Y)$ generated by $Y$, 
then the conditional expectation of $X$ given $Y$ is fine enough to find the averaging of $X$ when this event occurs.   
That is, $\E[X\Ind{A}] = \E[\E[X|Y]\Ind{A}]$.

\item If $X \in L^1$, then the conditional expectation $\E[X|Y] \in L^1$. 
\end{enumerate}
\end{rem}



\section{Conditional distribution given an event space} 
\begin{defn}%[Conditional probability given an event space] 
The \textbf{conditional probability %$P(A|\sG)$ 
of an event $A\in\sF$ given event space $\sG$} is 
defined as $P(A\mid \sG) \triangleq \E[\Ind{A}\mid\sG]$. 
\end{defn}
\begin{rem}
From the definition of conditional expectation, 
it follows that $P(A\mid \sG):\Omega\to[0,1]$ is a $\sG$ measurable random variable, such that 
$\E[\Ind{G}P(A|\sG)] = P(A\cap G)$ for all $G \in \sG$, 
and is uniquely defined up to sets of probability zero. 
\end{rem}
%\begin{defn}
%A conditional probability is called regular if $P(\cdot \mid \sG): \sG \to [0,1]$ is a probability measure on $(\Omega, \sF)$. 
%\end{defn}
\begin{shaded}
\begin{exmp} 
For the trivial sigma algebra   $\sG =\set{\emptyset ,\Omega}$, 
the conditional probability is the constant function
$P(A\mid \set{\emptyset ,\Omega})=P (A).$
\end{exmp}
\begin{exmp} 
If $A\in \sG$, then $P(A\mid\sG)=\Ind{A}$. 
\end{exmp}
\end{shaded}
\begin{defn}%[Conditional distribution given an event space]
%Consider a random variable $X:\Omega \to \R$ defined on probability space $(\Omega,\sF, P)$ such that $\E\abs{X} < \infty$ and a sub event space $\sG \subseteq \sF$. 
%Then, we can define 
The \textbf{conditional distribution of random variable $X$ given sub event space $\sG$} is defined as 
$
F_{X \mid \sG}(x)\triangleq P(A_X(x)\mid \sG)\text{ for all }x \in \R.
$
\end{defn}
\begin{rem} 
Recall that $F_{X\mid\sG}(x): \Omega\to[0,1]$ a random variable, for each $x\in\R$. 
Further, we observe that $F_{X\mid\sG}$ is monotone nondecreasing in $x \in \R$,  right continuous at all $x\in\R$, and has limits $\lim_{x\downarrow-\infty}F_{X\mid\sG}(x) = 0$ and $\lim_{x\uparrow\infty}F_{X\mid\sG}(x) = 1$. 
It follows that $F_{X\mid\sG}: \Omega \to [0,1]^\R$ is a random distribution. 
\end{rem}
\begin{thm} 
Let $g:\R\to\R$ be a Borel measurable function and $\sG$ be a sub-event space. 
Then, the conditional expectation $\E[g(X)\mid\sG] = \int_{x\in\R}g(x)dF_{X\mid\sG}(x).$
\end{thm}
\begin{proof}
It suffices to show this for simple random variables $X:\Omega\to\sX$. 
Since $g$ is Borel measurable, then $g(X)$ is a random variable. 
We will show that $\E[g(X)\mid\sG] = \sum_{x\in\sX}g(x)P_{X\mid\sG}(x)$ by showing that it satisfies three properties of conditional expectation.  
For part (i), we observe that from the definition of conditional probability $P_{X\mid\sG}(x)$ is a $\sG$-measurable random variable for all $x\in\sX$, and so is the linear combination $ \sum_{x\in\sX}g(x)P_{X\mid\sG}(x)$. 
For part (ii), we let $G \in \sG$. 
Then, it follows from the  linearity of expectation and the definition of conditional probability, that 
\EQ{
\E[\sum_{x\in\sX}g(x)P_{X\mid\sG}(x)\Ind{G}]
= \sum_{x\in\sX}g(x)\E[P_{X\mid\sG}(x)\Ind{G}] = \sum_{x\in\sX}g(x)\E[\Ind{E_x\cap G}]
= \E[X\Ind{G}].
}
For part (iii), it follows from the triangle inequality, the linearity of expectation, and the definition of conditional probability that
$
\E\abs{\sum_{x\in\sX}g(x)P_{X\mid\sG}(x)} \le \sum_{x\in\sX}\abs{g(x)} \E P_{X\mid\sG}(x)
= \sum_{x\in\sX}\abs{g(x)}P_X(x) = \E\abs{X} < \infty.
$
\end{proof}
\begin{rem} 
The conditional characteristic function is given by 
$\Phi_{X\mid\sG}(u) = \E[e^{juX}\mid\sG] = \int_{x\in\R}e^{jux}dF_{X\mid\sG}(x).$
\end{rem}
\begin{defn}%[Conditional distribution given a random variable]
%Consider two random variables $X, Y:\Omega\to\R$ defined on the same probability space $(\Omega, \sF, P)$ 
%such that $\E\abs{X} < \infty$. 
The \textbf{conditional distribution of random variable $X$ given random variable $Y$} is defined as $F_{X\mid Y}(x) \triangleq P(A_X(x)\mid\sigma(Y)]$ for all $x\in \R$.
\end{defn}
\begin{shaded}
\begin{exmp}[Conditional distribution given simple random variables] 
%For a simple random variable $Y: \Omega \to \sY \subseteq \R$ defined on the probability space $(\Omega, \sF, P)$, 
%we define fundamental events $E_y \triangleq Y^{-1}\set{y} \in \sF$ for all $y \in \sY$. 
%Then the sequence of events $E \triangleq (E_y \in \sF: y \in \sY)$ partitions the sample space, 
%and we can write the event space generated by random variable $Y$ as $\sigma(Y) = (\cup_{y\in I}E_y: I \subseteq \sY)$. 
Consider a random variable $X: \Omega \to \R$ and a simple random variable $Y:\Omega\to\sY$ defined on the same probability space. 
Since random variables $F_{X\mid Y}(x) = \E[\Ind{A_X(x)}\mid Y]$ are $\sigma(Y)$ measurable, 
they can can be written as $F_{X\mid Y}(x) = \sum_{y\in\sY}\beta_{x,y}\Ind{E_y}$ for some $\beta_x\in\R^\sY$ and $E_y = Y^{-1}\set{y}$ for all $y \in \sY$. 
Further, we have $\E[F_{X\mid Y}(x)\Ind{E_y}] = \E[\Ind{A_X(x)}\Ind{E_y}]$ for any $y \in \sY$, 
which implies that $\beta_{x,y} =\frac{P(A_X(x)\cap E_y)}{P_Y(y)} = F_{X\mid E_y}(x)$ for any $y \in \sY$. 
It follows that $F_{X\mid Y}$ is a $\sigma(Y)$ measurable simple random variable. 
%with probability mass function $$
\end{exmp}
\end{shaded}

\begin{shaded}
\begin{exmp}[Conditional expectation]
Consider a random experiment of a fair die being thrown and a random variable $X:\Omega\to\R$ taking the value of the outcome of the experiment. 
That is, for outcome space $\Omega = [6]$ and event space $\sF = \cP(\Omega)$, we have  $X(\omega) = \omega$ with $P_X(x) = 1/6$ for $x \in [6]$. 
Define another random variable $Y = \SetIn{X \le 3}$.
Then the conditional expectation of $X$ given $Y$ is a random variable given by
\EQ{
\E[X|Y] = \E[X\mid\set{Y =1}]\SetIn{Y = 1} + \E[X\mid \set{Y=1}]\SetIn{Y = 0}
= 2\Ind{Y^{-1}(1)} + 5\Ind{Y^{-1}(0)}.
}
Since $P\set{Y = 1} = P\set{Y = 0} = 0.5$, it follows that that $\E[\E[X|Y]] = \E[X]  = 3.5$.
\end{exmp}
\end{shaded}

\begin{shaded}
\begin{exmp}[Conditional distribution]
Consider the zero-mean Gaussian random variable $N:\Omega\to\R$ with variance $\sigma^2$, and another independent random variable $Y \in \set{-1,1}$ with PMF $(1-p, p)$ for some $p \in [0,1]$. 
Let $X = Y + N$, then the conditional distribution of $X$ given simple random variable $Y$ is 
\EQ{
F_{X|Y} = F_{X\mid Y^{-1}(-1)} \Ind{Y^{-1}(-1)} + F_{X\mid Y^{-1}(1)} \Ind{Y^{-1}(1)} ,
}
where $F_{X\mid Y^{-1}(\mu)}$ is $\int_{-\infty}^xe^{-\frac{(t-\mu)^2}{\sigma^2}}dt$. 
\end{exmp}
\end{shaded}

\begin{defn} When $X, Y$ are both continuous random variables, 
there exists a joint density $f_{X,Y}(x,y)$ for all $(x, y) \in \R^2$. 
For each $y \in \sY$ such that $f_Y(y) > 0$, we can define a function $f_{X\mid Y}: \R^2 \to \R_+$ such that   
\EQ{
f_{X\mid Y}(x,y) \triangleq \frac{f_{X,Y}(x,y)}{f_Y(y)}, \text{ for all } x \in \R.
}
\end{defn}
\begin{tcolorbox}
\begin{exerc}
For continuous random variables $X,Y$, 
show that the function $f_{X\mid Y^{-1}(y)}$ is a density of continuous random variable $X$ for each $y\in \R$. 
\end{exerc}
\end{tcolorbox}
\end{document}

\subsection{Conditioning on simple random variables}
Consider two random variables $X, Y$ defined on the same probability space $(\Omega, \sF, P)$.  
If the random variable $Y: \Omega \to \sY$ is simple, then the events $(E_y: y \in \sY)$ partition the sample space, 
where $E_y \triangleq Y^{-1}\set{y}$. 
Further the simple random variable $Y$ has a PMF $P_Y \triangleq (P(E_y): y \in \sY)$.   
If $y \in \sY$ such that $P_Y(y)>0$, then we can define a function $F_{X|E_y}: \R\to[0,1]$ 
such that
\EQN{
\label{eqn:ConditionalDistributionGivenEvent}
F_{X\given E_y}(x) \triangleq P\left(\set{X \le x}\given E_y\right) = \frac{P\set{X \le x, Y =y}}{P_Y(y)}, \text{ for all } x \in \R.  
}
\begin{shaded}
\begin{exerc}
For simple random variable $Y: \Omega \to \sY$, 
show that the function $F_{X\given E_y}$ conditioned on the event $E_y = \set{Y=y}$ is a distribution. 
%, and is a function of $y \in \sY$. 
\end{exerc}
\end{shaded}
\begin{defn}
Consider random variables $X,Y$ defined on the same probability space $(\Omega, \sF, P)$. 
If $Y: \Omega \to \sY$ is simple, 
then for all $y \in \sY$ such that $P_Y(y) > 0$, 
we can define event $E_y \triangleq Y^{-1}\set{y}$ and the \textbf{conditional distribution of $X$ given event $E_y$} as $F_{X \given E_y}: \R \to [0,1]$ defined in~\eqref{eqn:ConditionalDistributionGivenEvent}.
\end{defn}
\begin{defn}
%For a simple random variable $Y \in \sY$, 
%the distribution $F_{X\given E_y}$ is called the \textbf{conditional distribution of $X$ given event $E_y$}. 
%% and is a function of the realization $y \in \sY$. 
Consider random variables $X, Y$ defined on the same probability space $(\Omega, \sF, P)$. 
If $Y:\Omega\to\sY$ is simple such that $E_y \triangleq Y^{-1}\set{y}$ for all $y \in \sY$, 
then the \textbf{conditional distribution of $X$ given $Y$} is denoted by $F_{X|Y}: \Omega \to [0,1]^\R$ and defined as 
\EQ{
F_{X|Y}(\omega) \triangleq \sum_{y \in\sY}F_{X|E_y}\Ind{E_y}(\omega).
}
\end{defn}
\begin{rem}
We observe that $F_{X|\set{Y=y}} = F_{X|E_y}$. 
Hence, the conditional distribution $F_{X|Y}$ is a measurable function of the simple random variable $Y = \sum_{y\in\sY}y\Ind{E_y}$, 
and hence it is a random variable.
%such that
%\EQ{
%F_{X\given Y}: \omega \mapsto F_{X\given Y = Y(\omega)}.
%} 
%We can write the conditional distribution $F_{X|Y} = \sum_{y \in\sY}F_{X|E_y}\Ind{E_y}$.
\end{rem}

\subsection{Conditional densities}
When $X, Y$ are both continuous random variables, 
there exists a joint density $f_{X,Y}(x,y)$ for all $(x, y) \in \R^2$. 
For each $y \in \sY$ such that $f_Y(y) > 0$, we can define a function $f_{X\given Y = y}: \R \to \R_+$ such that   
\EQ{
f_{X\given Y = y}(x) \triangleq \frac{f_{X,Y}(x,y)}{f_Y(y)}, \text{ for all } x \in \R.
}
\begin{shaded}
\begin{exerc}
For continuous random variables $X,Y$, 
show that the function $f_{X\given Y= y}$ is a density of continuous random variable $X$ for each $y\in \R$. 
\end{exerc}
\end{shaded}
%\begin{defn}
%The \textbf{conditional density of $X$ given $Y$} for continuous random variables $X,Y$ is defined as a measurable function of the random variable $Y$, and hence it is random variable such that 
%\EQ{
%f_{X\given Y}: \omega \mapsto f_{X \given Y = Y(\omega)}. %\frac{f_{X,Y}(x,Y(\omega))}{f_Y(Y(\omega))}.
%}
%\end{defn}
%%\begin{defn}
%%For continuous random variables $X, Y$, 
%%the distribution $F_{X\given Y = y}(x)$ is called the \textbf{conditional distribution of $X$ given $Y=y$} and is a function of the realization $y \in \sY$. 
%%The \textbf{conditional distribution of $X$ given $Y$} is the family of conditional distributions of $X$ give $Y =y$ for all $y \in \sY$. 
%%That is, 
%%\EQ{
%%F_{X\given Y}(x) \triangleq \set{F_{X\given \set{Y = y}}(x): y \in \sY}.
%%} 
%%\end{defn}
%%The conditional density $f_{X\given Y}(x|y)$ is a function of $y \in \R$. 

\section{Conditional expectation in terms of conditional distribution}
Since we have defined the conditional distribution and densities, we can define the conditional expectation given an event as an integration with respect to the conditional distribution given that event.  
In the following, we will assume the random variables $X,Y$ are defined on the same probability space $(\Omega, \sF, P)$ and $\E\abs{X} < \infty$ such that $\E X$ exists and is finite. 
\subsection{Simple random variables}
When $Y:\Omega\to\sY$ is a simple random variable, 
we can define the conditional expectation of $X$ given the event $E_y = Y^{-1}\set{y}$ for $P_Y(y) > 0$ as 
\EQ{
\E[X\given E_y] \triangleq \int_{x \in \R}xdF_{X|E_y}(x) = \int_{(x,t)\in\R^2}x\frac{d_xF_{X,Y}(x,t)}{P_Y(y)}\SetIn{t=y} =\frac{\E[X\Ind{E_y}]}{P_Y(y)}.
}
%\begin{rem} 
%The conditional expectation of  $X$ given the event $\set{Y=y}$ is a function of realization $y \in \sY$ for $P_Y(y) > 0$.
%\end{rem}
\begin{defn}[Conditional expectation for conditioning on simple random variables] 
When $Y$ is a simple random variable, we can define the conditional expectation of $X$ given the random variable $Y$ is a measurable function of random variable $Y$, denoted by $\E[X\given Y]: \Omega \to \R$ such that 
\EQ{
\E[X\given Y]: \omega  \mapsto \int_{x \in \R}xdF_{X|\set{Y = Y(\omega)}}(x).
}
Hence, we can write
\EQ{
\E[X\given Y] = \sum_{y \in \sY}\E[X\given E_y]\Ind{E_y} =\sum_{y \in\sY}\frac{\E[X\Ind{E_y}]}{P_Y(y)}\Ind{E_y}.
}
\end{defn}
\begin{rem}
The random variable $\E[X\given Y]$ takes value $\E[X\given E_y]$ with probability $P_Y(y)$ for all $y \in \sY$. 
\end{rem}
\begin{lem} 
For simple random variable $Y: \Omega \to \sY$, 
the mean of random variable $\E[X \given Y]$ is $\E[X]$. 
\end{lem}
\begin{proof}
Since $\E[X|Y]$ is a function of the random variable $Y$, we have 
\EQ{
\E[\E[X\given Y]] = \sum_{y \in \sY}P_Y(y)\E[X\given E_y] = \sum_{y \in \sY}\E[X\Ind{E_y}]=\E[X].
}
%From the definition of conditional distribution of $X$ given the event $Y =y$, we get $F_{X \given Y = y}(x)P_Y(y) = P\set{X \le x, Y = y}$. 
%Interchanging derivative, summation, and integration from the application of dominated convergence theorem, and from the law of total probability, we get 
%\EQ{
%\E[\E[X\given Y]] = \int_{x \in \R}x dF_X(x) = \E[X]. 
%}
\end{proof}
%\begin{proof}
%Since $\E[X|Y]$ is a function of the random variable $Y$, we have 
%\EQ{
%\E[\E[X\given Y]] = \sum_{y \in \sY}P_Y(y)\E[X\given Y = y] = \sum_{y \in \sY}P_Y(y)\int_{x \in \R}xdF_{X\given Y =y}(x).
%}
%From the definition of conditional distribution of $X$ given the event $Y =y$, we get $F_{X \given Y = y}(x)P_Y(y) = P\set{X \le x, Y = y}$. 
%Interchanging derivative, summation, and integration from the application of dominated convergence theorem, and from the law of total probability, we get 
%\EQ{
%\E[\E[X\given Y]] = \int_{x \in \R}x dF_X(x) = \E[X]. 
%}
%\end{proof}

\subsection{Continuous random variables} 
When $X, Y$ are continuous random variables, we can define the conditional expectation of $X$ given $Y$ as a continuous random variable $\E[X\given Y]: \Omega \to \R$ such that 
\EQ{
\E[X\given Y]: \omega \mapsto \int_{x \in \R}xf_{X|Y = Y(\omega)}(x)dx,
}
with the density $f_Y(y)$ for all $y \in \R$.  
\begin{lem} 
For continuous random variables $X, Y$, 
the mean of random variable $\E[X \given Y]$ is $\E[X]$. 
\end{lem}
\begin{proof}
Since $\E[X|Y]$ is a function of the random variable $Y$ and its density is $f_Y(y)$, we get 
\EQ{
\E[\E[X\given Y]] = \int_{y \in \R}dy f_Y(y)\E[X\given Y = y] = \int_{y \in \R}dy f_Y(y)\int_{x \in \R}xf_{X\given Y =y}dx.
}
From the definition of conditional density of $X$ given $Y =y$, we get $f_{X \given Y = y}(x)f_Y(y) = f_{X,Y}(x,y)$. 
Interchanging integrations from Fubini's theorem, and from the law of total probability, we get 
\EQ{
\E[\E[X\given Y]] = \int_{x \in \R}x dx  \int_{y \in \R} f_{X,Y}(x,y)dy = \int_{x \in \R}xf_X(x) dx = \E[X]. 
}
\end{proof}

\section{Conditional expectation conditioned on an event space}
%\begin{defn}[Event space generated by a random variable] 
%Let $X: \Omega \to \R$ be a random variable on the probability space $(\Omega, \sF, P)$, 
%then the smallest event space generated by the events of the form $X^{-1}(-\infty, x]$ for $x \in \R$ is called the \textbf{event space generated} by this random variable $X$, 
%and denoted by $\sigma(X)$. 
%\end{defn}
%\begin{shaded}
%\begin{exmp}[Indicator function] 
%Let $A \in \sF$ be an event, then $X = \mathbbm{1}_A$ is a random variable and 
%\EQ{
%X^{-1}(-\infty, x] = \begin{cases}
%\Omega, &x \ge 1,\\
%A^c, & x \in [0, 1),\\
%\emptyset, &x < 0.
%\end{cases}
%}
%This implies that the smallest event space generated by this random variable is 
%$\sigma(X) = \set{\emptyset, A, A^c, \Omega}$. 
%\end{exmp}
%\end{shaded}
%\begin{shaded}
%\begin{exmp}[Simple random variables] 
%Let $X$ be a simple random variable, then $X = \sum_{x \in \sX}x\mathbbm{1}_{A_x}$ where $(A_x = X^{-1}\set{x} \in \sF: x \in \sX)$ is a finite partition of the sample space $\Omega$. 
%Without loss of generality, we can denote $\sX  = \set{x_1, \dots, x_n}$ where $x_1 \le \dots \le x_n$. 
%Then, 
%\EQ{
%X^{-1}(-\infty, x] = \begin{cases}
%\Omega, &x \ge x_n,\\
%\cup_{j=1}^iA_{x_i}, & x \in [x_i, x_{i+1}), i \in [n-1],\\
%\emptyset, &x < x_1.
%\end{cases}
%}
%Then the smallest event space generated by the simple random variable $X$ is $\set{\cup_{x \in S}A_x: S \subseteq \sX}$. 
%\end{exmp}
%\end{shaded}
\begin{defn}[Conditional expectation]
Consider two random variables $X,Y$ defined on the same probability space $(\Omega, \sF, P)$. 
\begin{enumerate}[(i)]
\item The \textbf{conditional distribution} of the random variable $X$ given a non-trivial event $B$ generated by the random variable $Y$ is 
\EQ{
F_{X \given B}(x) \triangleq \frac{P(\set{X \le x}\cap B)}{P(B)}. 
}
We can define $F_{X\given B} = 0$ for trivial events $B \in \sigma(Y)$ such that $P(B) = 0$. 

\item The \textbf{conditional expectation} of random variable $X$ given any event $B$ generated by random variable $Y$ is given by 
\EQ{
\E[X \given B] = \int_{x \in \R}xdF_{X\given B}.
}
We can define $\E[X\given B] = 0$ for trivial events $B \in \sigma(Y)$ such that $P(B) = 0$. 

%\item  The \textbf{conditional distribution} of the random variable $X$ given the random variable $Y$, 
%is a function of the random variable $Y$, and hence a random variable. 
%In particular, for each event $\set{Y \le y}$
%\EQ{
%F_{X\given Y}\SetIn{Y \le y} = F_{X \given \set{Y \le y}}\SetIn{Y \le y}  
%= \frac{F_{X,Y}(\cdot,y)}{F_Y(y)}\SetIn{Y \le y} ~\text{ for all }y \in \R.
%}
%That is, this conditional distribution is the collection $(F_{X \given B}: B \in \sigma(Y))$ of conditional distributions for each event $B$ generated by the random variable $Y$. 
%The probability of the event $B$ and the probability of the conditional distribution $F_{X \given B}\mathbbm{1}_B$ is $P(B)$. 
%
%\item The \textbf{conditional expectation} of random variable $X$ given the random variable $Y$, 
%is a function of the random variable $Y$, and hence a random variable. 
%In particular, for each event $\set{Y \le y}$, we have the event 
%\EQ{
%\E[X\given Y]\SetIn{Y \le y} = \E[X \given \set{Y \le y}]\SetIn{Y \le y}  
%= \SetIn{Y \le y} \int_{x \in \R}x dF_{X\given \set{Y \le y}}(x),\text{ for all }y \in \R.
%}
%That is, this conditional expectation is the collection of events $ (\E[X|B]: B \in \sigma(Y))$ corresponding to the  conditional expectation given the events generated by the random variable $Y$. 
%The probability of the event $B$ and the probability of the conditional expectation $\E[X\given B]\mathbb{1}_B$ is $P(B)$. 
%
%\item The mean of random variable $\E[X\given Y]\SetIn{Y \le y}$ for all $y \in \R$ is given by 
%\EQ{
%\E[\E[X\given Y]\SetIn{Y \le y}] 
%= \int_{t \in \R}dF_Y(t)\SetIn{Y \le y}  \int_{x \in \R}x dF_{X\given \set{Y \le y}}(x) 
%= \int_{(x,t) \in \R^2}x \SetIn{t \le y}dF_{X,Y}(x,t)
%= \E[X\SetIn{Y \le y}].
%}
%Similarly, for all events $B \in \sigma(Y)$, we have the event 
%\EQ{
%\E[\E[X\given Y]\mathbbm{1}_B] = \E[X\mathbbm{1}_B].
%}
\item The conditional expectation $\E[X\given Y]$ is a function of $Y$ and $\sigma(Y)$ measurable, 
and for all events $B \in \sigma(Y)$, we have 
\EQ{
\E[\E[X\given Y]\mathbbm{1}_B] = \E[X\mathbbm{1}_B]. 
}

\end{enumerate}
\end{defn}
%\begin{shaded}
%\begin{exmp}[Conditional expectation given indicator function] 
%Let $X$ be a random variable and $Y = \mathbbm{1}_B$ for some event $B \in \sF$.  
%\end{exmp}
%\end{shaded}

%\begin{defn}
%A random variable is independent of an event space $\sG$ if 
%\EQ{
%P(X^{-1}(-\infty, x]\cap  B) = F_X(x)P(B),\text{ for all }x \in \R, B \in \sG.
%}
%\end{defn}


\end{document}